% fixture: amsart/amsart-proofs
% oracle: pdflatex
% tier: T3
% engine: document
% covers: amsart: \qedhere in the display that ends a proof (\eqno\hbox{\qedsymbol}, no symbol at the end), gather* and align* (amsmath's, tags on the left: \eqnshift@ and \kern-\gdisplaywidth@\rlap{\boxz@}), \cite in amsart's \@citestyle (upright medium inside a theorem), thebibliography with [n] labels (no amsplain), \listparindent zero in \footnotesize's \@listi, \sloppy's \emergencystretch in the bibliography's font
% issue: #360
\documentclass[11pt]{amsart}
\newtheorem{theorem}{Theorem}
\begin{document}
\begin{theorem}
Every prime congruent to one modulo four is a sum of two squares \cite{HW,IK}, uniquely up to order and signs.
\end{theorem}
\begin{proof}
By the identity of Brahmagupta and Fibonacci it suffices to treat primes, and then
\[
(a^2+b^2)(c^2+d^2) = (ac-bd)^2 + (ad+bc)^2. \qedhere
\]
\end{proof}
A gathered display:
\begin{gather*}
x^2+y^2=z^2 \\
\sum_{i=1}^n i=\frac{n(n+1)}{2}
\end{gather*}
and an aligned one:
\begin{align*}
a &= \int_0^1 f(x)\,dx \\
  &= b.
\end{align*}
\begin{thebibliography}{9}
\bibitem{HW} G.~H. Hardy and E.~M. Wright, \emph{An introduction to the theory of numbers}, Oxford, 2008.
\bibitem{IK} H.~Iwaniec and E.~Kowalski, \emph{Analytic number theory}, AMS, 2004.

A second paragraph in the same item, not indented in the bibliography's footnote size.
\bibitem{L} See \texttt{verylongidentifierthatcannotbebrokenanywhereatall} and \texttt{anotherverylongidentifierthatcannotbebroken}.
\end{thebibliography}
\end{document}
